\documentclass[11pt,a4paper]{article}
\usepackage[textwidth=182mm,top=20mm,bottom=22mm]{geometry}
\usepackage{amsmath,amssymb,newtxmath}
\usepackage[no-math]{fontspec}
\setmainfont{Times New Roman}
\usepackage{xeCJK}
\setCJKmainfont[AutoFakeBold=2]{SimSun}
\usepackage{graphicx,booktabs,array,tabularx,seqsplit,xurl,hyperref}
\hypersetup{hidelinks,unicode=true,pdftitle={研究论文模板},pdfauthor={作者姓名}}
\setlength{\emergencystretch}{2em}
\raggedbottom
\renewcommand{\arraystretch}{1.15}
\usepackage{caption}
\captionsetup[figure]{font=small,justification=centering,singlelinecheck=false}
\linespread{1.16}
\setlength{\parindent}{2em}
\setlength{\parskip}{3pt plus 1pt}
\renewcommand{\topfraction}{0.92}
\renewcommand{\textfraction}{0.08}
\renewcommand{\floatpagefraction}{0.75}
\renewcommand{\figurename}{图}
\renewcommand{\tablename}{表}
\renewcommand{\abstractname}{摘要}
\begin{document}
\title{研究论文模板}
\author{作者姓名\\单位名称}
\date{}
\maketitle
\begin{abstract}
概述研究问题、方法、主要结果和适用范围。
\end{abstract}
\noindent\textbf{关键词：}模型；验证\par
\section{方法}
\hypertarget{chapter-01}{}
使用 SI 单位和 RMS 相量。圆线解的特殊函数定义见文献 [1]。

\begin{equation}\delta=\frac{1}{\sqrt{\pi f\mu\sigma}}\label{eq:skin}\end{equation}
\begin{table}[!htbp]
\centering
\caption{符号及单位}
\begingroup\footnotesize
\begin{tabularx}{\textwidth}{>{\raggedright\arraybackslash}X>{\raggedright\arraybackslash}X}
\toprule
Symbol & Unit\\
\midrule
f & Hz\\
σ & S/m\\
\bottomrule
\end{tabularx}\endgroup
\end{table}
\section{结果与讨论}
\hypertarget{chapter-02}{}
将图件替换为可追溯的研究结果，说明比较条件、误差与结论。

\renewcommand{\refname}{参考文献}
\begin{figure}[!htbp]\centering\fbox{\rule{0pt}{22mm}\rule{.8\linewidth}{0pt}}\caption{图件位置示例；请替换为实际结果。}\label{fig:example}\end{figure}
\begin{thebibliography}{99}
\bibitem[1]{BESSEL} NIST. NIST Digital Library of Mathematical Functions, Modified Bessel Functions. 2026. \url{https://dlmf.nist.gov/10.25}
\end{thebibliography}
\end{document}
